\documentclass[12pt]{article}
\usepackage{graphicx}
\usepackage[T1]{fontenc}
\usepackage{tocloft}
\usepackage{todonotes}
\usepackage{subcaption}
\usepackage{caption}
\usepackage{hyperref}
\usepackage{booktabs}
\usepackage{listings}
\usepackage{wrapfig}
\usepackage{pdfpages}
\usepackage{pdflscape}
\usepackage{scrlayer}
\usepackage{scrhack}
\begin{document}
\begin{titlepage}
    \centering
    {\scshape\large AY 2020/2021 \par}
    \vfill
    \includegraphics[width=100pt]{assets/logo-polimi-new}\par\vspace{1cm}
    {\scshape\LARGE Politecnico di Milano \par}
    \vspace{1.5cm}
    {\huge\bfseries RASD\@: Requirement Analysis
        and Specification Document \par}
    \vspace{2cm}
    {\Large {Alice Piemonti\quad Luca Pirovano\quad Nicolò Sonnino}\par}
    \vfill
    {\large Professor\par
        Matteo \textsc{Rossi}}
    \vfill
    {\large \today \par}
\end{titlepage}
\tableofcontents
\newpage
\section{Introduction}
CLup (Customers Line-up) is an easy-to-use application which intent is to help grocery shopping to face the big challenges that arise during this tough period.\\

In fact, due to the recent worldwide spread of SARS-CoV-2 (COVID-19), many countries are imposing lockdowns and strict regulations about social distancing such as: the closure of restaurants in the evening, limitations on public transports, curfews, etc.

In particular, supermarkets are imposed to restrict accesses to their stores, in order to avoid having too much people inside, which stands at the basis of social distancing. Furthermore, the staggered access needed by stores lead to the same problem in the outside, such as long queues and crowds.

Grocery stores need an application which main intent is to manage in an efficient way customers arrival on the outside of the supermarket, and to count the access of people to the inside. The advantage is to regulate the influx of people near the stores, according to the vigent strict rules, and to avoid long hours of waiting, saving people from line up.

\subsection{Purpose}
The aim of the product is to avoid gatherings outside and inside grocery stores,  improving the safety of the customers.

This is achieved through monitoring accesses to the buildings, managing time slots for visits and optimizing people flows inside the stores.\\

The application should provide two types of accesses for customers (as normal users) and stores' attendants (as special users).

Customers will have the possibility to line up in a virtual queue, so that they can wait from a close and safe building until their number is called. In an accettable time, the application will inform the user when his number is about to be called, so that he can reach the store in the right time.

In addition, users can book a visit to the supermarket for a different time or day. Hence, the customer indicates the slot preferred and, if it is available the application register the reservation, otherwise a list of alternatives are displayed, such as the possibility to book in different slot, or to chose differnt stores available for that time/day.

The application can also register in advance the duration of the visit of a customer and a list of categories that the customer indends to buy, so that the system can suggest the best slot and plan in a finer way the visits, in order to guarantee a correct distance through the aisles inside the store. In fact, the application will be able to balance the presence of people in all the areas of the supermarket, as well as the flows of visitors throughout the day.

The customer will be able to activate an additional functionality: the application will send a notification of available slots in a certain day/time range, in order to leave to the user the minimum effort.

Attendants will have the possibility to scan a QR code at the entance of the sotre, in order to verify the correctness of the arrival of the customer and, in the meanwhile, monitor the number of entrances.

In addition, attendants should have access to a proper area of the application that permits to behave as a proxy and hand out tickets on spot in order to guarantee the access to those people who doesn't have the application.

The application will be operable freely, widely available and very intuitive, because the range of users (i.e. people who need to go to the grocery store) extends to the entire population.

The userbase is expected to be both people with an Internet access and ones without it, from young people to elderly, thanks to the possibility of attendants to act as a proxy.

\subsection{Scope}
The product shall be called CLup and will let users to plan their shopping session in two different ways:
\begin{itemize}
    \item {\textbf{ASAP}: the user will claim the first available ticket and receive an estimated queue time.}
    \item {\textbf{Reservation}: the user will choose a time slot from a list of available ones, in order to book his visit to the structure.}
\end{itemize}

Every customer can choose one of these modes \textbf{remotely} via an official app or through a web browser, or \textbf{in presence} by asking to a staff member, who will act as an intermediate between the customer and the system.

When a customer makes a reservation, the system allows him to choose the duration of his visit and insert a list of possible purchases, in order to optimize his stay.\\

In addition, the user can change time slot/store based on system's suggestions and enable periodically notifications of available slots in a day/time range.

\subsection{Goals}
The main objectives of our system are the following:
\begin{itemize}
    \item {\textbf{G1: Allow customers to retrieve a queue number}\\This is the main feature of the application, through which customers are forced not staying outside the structure.}
    \item {\textbf{G2: Allow customers to know an estimated queue time}\\Through an appropriate estimation of each customer's permanence time, the user is given an estimated queue time to let going to the supermarket when needed.}
    \item {\textbf{G3: Guarantee that each queue number is unique}\\There is no way that a queue number could be assigned to more than one customer.}
    \item {\textbf{G4: Allow customers to generate a QR Code}\\A QR Code would let store managers to monitor entrances, scanning a customer's code upon entering.}
    \item {\textbf{G5: Allow shops to offer the remote queue function}\\Each shop can register to the service and offer its customer to Line-Up from home.}
    \item {\textbf{G6: Allow shops to generate tickets on the spot}\\If someone does not have access to the required technology, they can still take advantage of the system by getting their queue number directly at the store.}
    \item {\textbf{G7: Allow customers to "book a visit"}\\Customers can book a slot providing the expected duration of the visit and, also, they can provide a list of categories in which the items they want to buy belong. In this case, they will occupy different slots and so there is the possibility of optimizing visits.}
    \item {\textbf{G8: Let the system infer a visit duration}\\For long term customer, system should provide an estimation of visit duration, relying on their previous shopping sessions.}
    \item {\textbf{G9: Allow customers to receive a suggestion of alternative slots}\\The aim of this goal is to balance out the number of people in the store. The suggestion could be on the same shop chain or also on different chains.}
    \item {\textbf{G10: Allow customers to receive notifications on free slots}\\Customers could enable notifications in order to receive the first available slot and, if needed, to book it.}
\end{itemize}

\newpage
\section{Overall Description}
\newpage
\section{Specific Requirements}
\newpage
\section{Formal Analysis}
\newpage
\section{Effort spent}
\end{document}